%% Created by Maple 2022.1, Windows 10
%% Source Worksheet: lab_sols10
%% Generated: Wed May 08 10:35:20 BST 2024
\documentclass{article}
\usepackage{amssymb}
\usepackage{graphicx}
\usepackage{hyperref}
\usepackage{listings}
\usepackage{mathtools}
\usepackage{maple}
\usepackage[utf8]{inputenc}
\usepackage[svgnames]{xcolor}
\usepackage{amsmath}
\usepackage{breqn}
\usepackage{textcomp}
\begin{document}
\lstset{basicstyle=\ttfamily,breaklines=true,columns=flexible}
\pagestyle{empty}
\DefineParaStyle{Maple Bullet Item}
\DefineParaStyle{Maple Heading 1}
\DefineParaStyle{Maple Warning}
\DefineParaStyle{Maple Heading 4}
\DefineParaStyle{Maple Heading 2}
\DefineParaStyle{Maple Heading 3}
\DefineParaStyle{Maple Dash Item}
\DefineParaStyle{Maple Error}
\DefineParaStyle{Maple Title}
\DefineParaStyle{Maple Ordered List 1}
\DefineParaStyle{Maple Text Output}
\DefineParaStyle{Maple Ordered List 2}
\DefineParaStyle{Maple Ordered List 3}
\DefineParaStyle{Maple Normal}
\DefineParaStyle{Maple Ordered List 4}
\DefineParaStyle{Maple Ordered List 5}
\DefineCharStyle{Maple 2D Output}
\DefineCharStyle{Maple 2D Input}
\DefineCharStyle{Maple Maple Input}
\DefineCharStyle{Maple 2D Math}
\DefineCharStyle{Maple Hyperlink}
\begin{center}
\begin{Maple Normal}
\underline{\textbf{Taylor series}}
\end{Maple Normal}
\end{center}
\begin{Maple Heading 2}
\textbf{Exercise 1.1}
\end{Maple Heading 2}
\begin{Maple Normal}
(a)
\end{Maple Normal}
\begin{lstlisting}
> y := ln(sqrt((1+x)/(1-x)));
\end{lstlisting}
\begin{lstlisting}
> simplify(diff(y,x$5));
\end{lstlisting}
\begin{lstlisting}
> subs(x=0,%);
\end{lstlisting}
\begin{lstlisting}
> %/5!;
\end{lstlisting}
\begin{Maple Normal}
(b)
\end{Maple Normal}
\begin{lstlisting}
> simplify(subs(x=0,y));
\end{lstlisting}
\begin{lstlisting}
> subs(x=0,simplify(diff(y,x$1)))/1!;
\end{lstlisting}
\begin{lstlisting}
> subs(x=0,simplify(diff(y,x$2)))/2!;
\end{lstlisting}
\begin{lstlisting}
> subs(x=0,simplify(diff(y,x$3)))/3!;
\end{lstlisting}
\begin{lstlisting}
> subs(x=0,simplify(diff(y,x$4)))/4!;
\end{lstlisting}
\begin{lstlisting}
> subs(x=0,simplify(diff(y,x$5)))/5!;
\end{lstlisting}
\begin{Maple Normal}
(c)
\end{Maple Normal}
\begin{lstlisting}
> a := (n) -> subs(x=0,diff(y,x$n))/n!;
\end{lstlisting}
\begin{lstlisting}
> seq(a(n),n=1..5);
\end{lstlisting}
\begin{Maple Normal}
(d)
\end{Maple Normal}
\begin{lstlisting}
> add(a(k)*x^k,k=1..12);
\end{lstlisting}
\begin{lstlisting}
> series(y,x=0,13);
\end{lstlisting}
\begin{Maple Normal}
(e)
\end{Maple Normal}
\begin{Maple Normal}
The obvious guess is that 
$y$ is the sum of  
$\frac{x^{n}}{n}$ for all odd integers 
$n$, or in other words 
$y =\moverset{\infty}{\munderset{k =0}{\sum}}\frac{x^{2 k +1}}{2 k +1}$.\

This can be checked as follows:
\end{Maple Normal}
\begin{lstlisting}
> sum(x^(2*k+1)/(2*k+1),k=0..infinity);
\end{lstlisting}
\begin{lstlisting}
> simplify(%-y,symbolic);
\end{lstlisting}
\begin{lstlisting}
> 
\end{lstlisting}
\begin{Maple Normal}
\_\_\_\_\_\_\_\_\_\_\_\_\_\_\_\_\_\_\_\_\_\_\_\_\_\_\_\_\_\_\_\_\_\_\_\_\_\_\_\_\_\_\_\_\_\_\_\_\_\_\_\_\_\_\_\_\_\_\_\_\_\_\_\_\_\_\_\_\_\_\_\_\_\_\_\_\_\_\_\_\_
\end{Maple Normal}
\begin{Maple Heading 2}
\textbf{Exercise 1.2}
\end{Maple Heading 2}
\begin{Maple Normal}
(a)
\end{Maple Normal}
\begin{lstlisting}
> s := convert(series(sin(x),x=0,12),polynom);
\end{lstlisting}
\begin{Maple Normal}
(b)
\end{Maple Normal}
\begin{lstlisting}
> plot([s,sin(x)],x=-8..8,-10..10);
\end{lstlisting}
\begin{Maple Normal}
(c)
\end{Maple Normal}
\begin{lstlisting}
> t := (n) -> convert(series(sin(x),x=0,n),polynom);
\end{lstlisting}
\begin{lstlisting}
> plot([t(14),sin(x)],x=-8..8,-10..10);
\end{lstlisting}
\begin{Maple Normal}
The graph of 
$t (n)$  (for 
$x$ in 
$[-8,8]$) is visually indistinguishable from that of 
$\sin (x)$ when 
$20\le n$.
\end{Maple Normal}
\begin{lstlisting}
> plot([t(20),sin(x)],x=-8..8,-10..10);
\end{lstlisting}
\begin{Maple Normal}
However, the two functions diverge sharply for 
$x$ outside this range:
\end{Maple Normal}
\begin{lstlisting}
> plot([t(20),sin(x)],x=-10..10,-10..10);
\end{lstlisting}
\begin{lstlisting}
> 
\end{lstlisting}
\begin{Maple Normal}
(d)
\end{Maple Normal}
\begin{lstlisting}
> plot([sin(x),seq(t(4*n+2),n=0..10)],x=0..20,-2..10);
\end{lstlisting}
\begin{lstlisting}
> plot([sin(x),seq(t(4*n+4),n=0..10)],x=0..20,-10..2);
\end{lstlisting}
\begin{Maple Normal}
We do not bother with 
$t (n)$ for odd 
$n$, because 
$t (n)=t (n -1)$ for odd 
$n$.
\end{Maple Normal}
\begin{Maple Normal}

\end{Maple Normal}
\begin{Maple Normal}
(e)
\end{Maple Normal}
\begin{lstlisting}
> r := (n) -> convert(series(cos(x),x=0,n),polynom);
\end{lstlisting}
\begin{lstlisting}
> expand(t(10)^2+r(10)^2);
\end{lstlisting}
\begin{Maple Normal}

\end{Maple Normal}
\begin{Maple Normal}
Note that 
$t (10)$ is an approximation to 
$\sin (x)$ and 
$r (10)$ is an approximation to 
$\cos (x)$, so 
$t (10)^{2}+r (10)^{2}$
\end{Maple Normal}
\begin{Maple Normal}
should be an approximation to 
$\sin (x)^{2}+\cos (x)^{2}=1$.  In fact we see that  
$t (10)^{2}+r (10)^{2}$ is 1 plus a sum
\end{Maple Normal}
\begin{Maple Normal}
of terms like 
$\frac{x^{k}}{N}$ where 
$k$ is at least 
$10$ and 
$N$ is at least 
$1000000$.  If 
$x$ is not too big, then all these extra 
\end{Maple Normal}
\begin{Maple Normal}
terms will be very small, so 
$t (10)^{2}+r (10)^{2}$ will indeed be close to 
$1$.  To see this graphically, we plot
\end{Maple Normal}
\begin{Maple Normal}

$1-t (10)^{2}-r (10)^{2}$ :
\end{Maple Normal}
\begin{lstlisting}
> plot(1-t(10)^2-r(10)^2,x=-5..5,-1..1);
\end{lstlisting}
\begin{Maple Normal}
The error is very small for 
$x$ between 
$-3$ and 
$3$, but it grows very big when 
$x <-5$ or 
$5<x$.
\end{Maple Normal}
\begin{Maple Normal}
\_\_\_\_\_\_\_\_\_\_\_\_\_\_\_\_\_\_\_\_\_\_\_\_\_\_\_\_\_\_\_\_\_\_\_\_\_\_\_\_\_\_\_\_\_\_\_\_\_\_\_\_\_\_\_\_\_\_\_\_\_\_\_\_\_\_\_\_\_\_\_\_\_\_\_\_\_\_\_\_\_
\end{Maple Normal}
\begin{Maple Heading 2}
\textbf{Exercise 1.3}
\end{Maple Heading 2}
\begin{lstlisting}
> series(x*(1+x)/(1-x)^3,x=0,11);
\end{lstlisting}
\begin{Maple Normal}

\end{Maple Normal}
\begin{Maple Normal}
Note that the coefficient of 
$x^{3}$ is 
$9=3^{2}$, the coefficient of 
$x^{4}$ is 
$16=4^{2}$ and so on.  It should be 
\end{Maple Normal}
\begin{Maple Normal}
clear from this that the series is 
$\moverset{\infty}{\munderset{k =1}{\sum}}k^{2} x^{k}$.  We can ask Maple to confirm this as follows:
\end{Maple Normal}
\begin{Maple Normal}

\end{Maple Normal}
\begin{lstlisting}
> sum(k^2*x^k,k=1..infinity);
\end{lstlisting}
\begin{Maple Normal}

\end{Maple Normal}
\begin{Maple Normal}
This is the same as 
$\frac{x (1+x)}{(1-x)^{3}}$,  rearranged slightly.
\end{Maple Normal}
\begin{Maple Normal}
\_\_\_\_\_\_\_\_\_\_\_\_\_\_\_\_\_\_\_\_\_\_\_\_\_\_\_\_\_\_\_\_\_\_\_\_\_\_\_\_\_\_\_\_\_\_\_\_\_\_\_\_\_\_\_\_\_\_\_\_\_\_\_\_\_\_\_\_\_\_\_\_\_\_\_\_\_\_\_\_\_
\end{Maple Normal}
\begin{Maple Heading 2}
\textbf{Exercise 1.4}
\end{Maple Heading 2}
\begin{Maple Normal}

\end{Maple Normal}
\begin{lstlisting}
> series(sin(x),x=Pi/2,12);
\end{lstlisting}
\begin{lstlisting}
> series(cos(x),x=0,12);
\end{lstlisting}
\begin{Maple Normal}

\end{Maple Normal}
\begin{Maple Normal}
If we call the first series 
$f (x)$ and the second series 
$g (x)$, then the relationship is that 
$f (x)=g (x -\frac{\mathrm{Pi}}{2})$.
\end{Maple Normal}
\begin{Maple Normal}
This is reasonable, because 
$f (x)$ approximates 
$\sin (x)$ and 
$g (x)$ approximates 
$\cos (x)$, and
\end{Maple Normal}
\begin{Maple Normal}

$\sin (x)=\cos (x -\frac{\mathrm{Pi}}{2})$.
\end{Maple Normal}
\begin{Maple Normal}
\_\_\_\_\_\_\_\_\_\_\_\_\_\_\_\_\_\_\_\_\_\_\_\_\_\_\_\_\_\_\_\_\_\_\_\_\_\_\_\_\_\_\_\_\_\_\_\_\_\_\_\_\_\_\_\_\_\_\_\_\_\_\_\_\_\_\_\_\_\_\_\_\_\_\_\_\_\_\_\_\_
\end{Maple Normal}
\begin{Maple Heading 2}
\textbf{Exercise 2.1}
\end{Maple Heading 2}
\begin{lstlisting}
> evalf[20](cos(ln(Pi+20)));
\end{lstlisting}
\begin{Maple Normal}
No one seems to have a good explanation for why this is so close to -1.  It is probably just a coincidence.
\end{Maple Normal}
\begin{Maple Normal}
\_\_\_\_\_\_\_\_\_\_\_\_\_\_\_\_\_\_\_\_\_\_\_\_\_\_\_\_\_\_\_\_\_\_\_\_\_\_\_\_\_\_\_\_\_\_\_\_\_\_\_\_\_\_\_\_\_\_\_\_\_\_\_\_\_\_\_\_\_\_\_\_\_\_\_\_\_\_\_\_\_
\end{Maple Normal}
\begin{Maple Heading 2}
\textbf{Exercise 2.2}
\end{Maple Heading 2}
\begin{lstlisting}
> f := (x) -> x^2-29/16;
\end{lstlisting}
\begin{lstlisting}
> f(0);
\end{lstlisting}
\begin{lstlisting}
> f(f(f(-1/4)));
\end{lstlisting}
\begin{Maple Normal}
We see that 
$-\frac{1}{4}$ is sent to itself by the function 
$f (f (f (x)))$.  There is a big theory about points fixed under
\end{Maple Normal}
\begin{Maple Normal}
iterated application of functions like 
$f (x)$, which you can learn in the level 3 Chaos course.
\end{Maple Normal}
\begin{Maple Normal}
\_\_\_\_\_\_\_\_\_\_\_\_\_\_\_\_\_\_\_\_\_\_\_\_\_\_\_\_\_\_\_\_\_\_\_\_\_\_\_\_\_\_\_\_\_\_\_\_\_\_\_\_\_\_\_\_\_\_\_\_\_\_\_\_\_\_\_\_\_\_\_\_\_\_\_\_\_\_\_\_\_
\end{Maple Normal}
\begin{Maple Heading 2}
\textbf{Exercise 2.3}
\end{Maple Heading 2}
\begin{lstlisting}
> a := (1+x)^5-3*(1+x)^4+5*(1+x)^3-3*(1+x)^2+3*(1+x)+3;
\end{lstlisting}
\begin{lstlisting}
> b := (7*x^2-6*x-x^8)/(x-1)^2;
\end{lstlisting}
\begin{lstlisting}
> simplify(a);
\end{lstlisting}
\begin{lstlisting}
> simplify(b);
\end{lstlisting}
\begin{Maple Normal}

\end{Maple Normal}
\begin{Maple Normal}
It is clear from this that 
$b =-x a$.
\end{Maple Normal}
\begin{Maple Normal}
\_\_\_\_\_\_\_\_\_\_\_\_\_\_\_\_\_\_\_\_\_\_\_\_\_\_\_\_\_\_\_\_\_\_\_\_\_\_\_\_\_\_\_\_\_\_\_\_\_\_\_\_\_\_\_\_\_\_\_\_\_\_\_\_\_\_\_\_\_\_\_\_\_\_\_\_\_\_\_\_\_
\end{Maple Normal}
\begin{Maple Heading 2}
\textbf{Exercise 2.4}
\end{Maple Heading 2}
\begin{lstlisting}
> y := ((x^12-1)*(x^2-1)/((x^6-1)*(x^4-1)))^(10);
\end{lstlisting}
\begin{lstlisting}
> simplify(y);
\end{lstlisting}
\begin{lstlisting}
> coeff(simplify(y),x^6);
\end{lstlisting}
\begin{lstlisting}
> 
\end{lstlisting}
\begin{Maple Normal}
\_\_\_\_\_\_\_\_\_\_\_\_\_\_\_\_\_\_\_\_\_\_\_\_\_\_\_\_\_\_\_\_\_\_\_\_\_\_\_\_\_\_\_\_\_\_\_\_\_\_\_\_\_\_\_\_\_\_\_\_\_\_\_\_\_\_\_\_\_\_\_\_\_\_\_\_\_\_\_\_\_
\end{Maple Normal}
\begin{Maple Heading 2}
\textbf{Exercise 2.5}
\end{Maple Heading 2}
\begin{lstlisting}
> restart;
\end{lstlisting}
\begin{lstlisting}
> solve(\{
 x^2 + y^2 + z^2 = 9,
 (x-1)^2 + (y-1)^2 + (z-1)^2 = 2,
 4*x^2+y*z = 2*x*y+2*x*z
\},\{x,y,z\});
\end{lstlisting}
\begin{Maple Normal}

\end{Maple Normal}
\begin{Maple Normal}
We see that the only integer solution is 
$(x ,y ,z)=(1,2,2)$.  
\end{Maple Normal}
\begin{Maple Normal}
\_\_\_\_\_\_\_\_\_\_\_\_\_\_\_\_\_\_\_\_\_\_\_\_\_\_\_\_\_\_\_\_\_\_\_\_\_\_\_\_\_\_\_\_\_\_\_\_\_\_\_\_\_\_\_\_\_\_\_\_\_\_\_\_\_\_\_\_\_\_\_\_\_\_\_\_\_\_\_\_\_
\end{Maple Normal}
\begin{Maple Heading 2}
\textbf{Exercise 2.6}
\end{Maple Heading 2}
\begin{lstlisting}
> _EnvAllSolutions := true;
\end{lstlisting}
\begin{lstlisting}
> solve(sin(theta)^2=3*cos(theta)^2,\{theta\});
\end{lstlisting}
\begin{lstlisting}
> plot(sin(Pi*x)^2-3*cos(Pi*x)^2,x=-2..2);
\end{lstlisting}
\begin{Maple Normal}

\end{Maple Normal}
\begin{Maple Normal}
The solutions are 
$\theta =(n +\frac{1}{3}) \mathrm{Pi}$ and 
$\theta =(n -\frac{1}{3}) \mathrm{Pi}$ for integers 
$n$.
\end{Maple Normal}
\begin{Maple Normal}
\_\_\_\_\_\_\_\_\_\_\_\_\_\_\_\_\_\_\_\_\_\_\_\_\_\_\_\_\_\_\_\_\_\_\_\_\_\_\_\_\_\_\_\_\_\_\_\_\_\_\_\_\_\_\_\_\_\_\_\_\_\_\_\_\_\_\_\_\_\_\_\_\_\_\_\_\_\_\_\_\_
\end{Maple Normal}
\begin{Maple Heading 2}
\textbf{Exercise 2.7}
\end{Maple Heading 2}
\begin{lstlisting}
> _EnvExplicit := true;
\end{lstlisting}
\begin{lstlisting}
> sols := solve(\{x^2-y^2=1,2*x*y=1\},\{x,y\});
\end{lstlisting}
\begin{lstlisting}
> evalf(subs(sols[1],[x,y]));
\end{lstlisting}
\begin{lstlisting}
> evalf(subs(sols[2],[x,y]));
\end{lstlisting}
\begin{lstlisting}
> evalf(subs(sols[3],[x,y]));
\end{lstlisting}
\begin{lstlisting}
> evalf(subs(sols[4],[x,y]));
\end{lstlisting}
\begin{Maple Normal}
Only the first of these is a pair of positive real numbers.  Thus, the solution we want is as follows:
\end{Maple Normal}
\begin{lstlisting}
> sols[1];
\end{lstlisting}
\begin{lstlisting}
> simplify(%);
\end{lstlisting}
\begin{lstlisting}
> 
\end{lstlisting}
\begin{Maple Normal}
\_\_\_\_\_\_\_\_\_\_\_\_\_\_\_\_\_\_\_\_\_\_\_\_\_\_\_\_\_\_\_\_\_\_\_\_\_\_\_\_\_\_\_\_\_\_\_\_\_\_\_\_\_\_\_\_\_\_\_\_\_\_\_\_\_\_\_\_\_\_\_\_\_\_\_\_\_\_\_\_\_
\end{Maple Normal}
\begin{Maple Heading 2}
\textbf{Exercise 2.8}
\end{Maple Heading 2}
\begin{lstlisting}
> fsolve(x=log(x+20),x=3);
\end{lstlisting}
\begin{lstlisting}
> 
\end{lstlisting}
\begin{lstlisting}
> 
\end{lstlisting}
\begin{Maple Normal}
This is related to Exercise 2.1.  If the answer there were exactly -1, then the answer here would be 
$\mathrm{Pi}$.
\end{Maple Normal}
\begin{Maple Normal}
  
\end{Maple Normal}
\end{document}
